\section{Introduction}
\label{sec:introduction}

The Sieve Sequence represents an infinite accepted-value stream using one
finite cyclic gap list. That representation makes complete-period
identities exact, but the twin-prime application asks for a surviving
2-gap start in the eligible square-safe window of one future head. The
distinction between those scopes is the organizing principle of this
article.

Version 1 established the copy-or-merge and complete-period/local-window
boundary. This version retains that theorem spine and adds the signed
localization results that emerged from the later quadratic investigation.
The article develops the following properties:

\begin{enumerate}
  \item exact complete-period 2-gap count --- Section~3.1;
  \item exact copy-index filter frequency --- Section~3.2;
  \item exact batched survival --- Section~3.3;
  \item exact complete-period \texttt{(2,4,2)} cluster count --- Section~3.4;
  \item rotation invariance of cyclic counts --- Section~3.5;
  \item stable global absence of 2-gaps --- Section~3.6;
  \item square-safe certification --- Section~4.1;
  \item post-filter-3 isolation --- Section~4.2;
  \item exact accepted strikes and the sharp one-transition threshold ---
    Sections~4.3--4.4;
  \item why the capacity envelope is exhausted --- Section~6;
  \item exact weighted deletion conservation and its terminal quadratic
    corollary --- Sections~5.1--5.2;
  \item the exact filter-$7$ interval saving --- Section~7;
  \item the live frontier (accepted-boundary discrepancy and
    residue-collision energy estimates) --- Section~8;
  \item the copy-block residue-energy bridge --- Section~9;
  \item the classified routes and the live program --- Section~10; and
  \item the fixed-seed scale conflict (Section~10.1) and the Type-II
    barrier (Section~11.1).
\end{enumerate}

The final sections separate the proved complete-block control from the
open partial-boundary and cross-layer estimates. No complete-period
theorem is used as a short-window lower bound.
